% !TEX root = ../../Main.tex
\section{Article Research: Summarising Different Implementations}

Next, eighteen articles with working implementations were read in full. Their features were recorded in two tables that cover the same articles. \Cref{Tables:BasicRelatedWork} holds the market, the data and the algorithm used. \Cref{Tables:AdvancedRelatedWork} holds the design choices inside each agent. Below, the articles are grouped by the market they trade.


\subsection{Forex Market Implementations}

Three of the eighteen articles trade currencies. All three limit the agent to a discrete choice between buying, holding and selling.

Carapuço et al. \cite{carapuco_reinforcement_2018} built a Deep Q-Network over a fully connected network and used it on raw tick data from 2010 to 2017. The agent could act only after a fixed number of ticks had passed. This gave about one decision every two hours. It outperformed its benchmark by a wide margin. It is also the only study reviewed here that works from tick data and not from bars. Huang \cite{huang_financial_2018} used a Deep Recurrent Q-Network on 15-minute bars from 2012 to 2017. This model puts a long short-term memory network under the Q-function. It outperformed its benchmark with explicit and implicit trading costs charged. This makes it one of the few results in the set that still hold when costs are modelled. Rundo \cite{rundo_deep_2019} used the same kind of model on one-minute bars from 2004 to 2018, and also reported that it outperformed its benchmark. However, its accuracy was high enough to suggest look-ahead bias, and its maximum drawdown was large. Taken together, the two results suggest that adding recurrence is not enough on its own to make very short timeframes workable.

\subsection{Stock Market Implementations}

Most of the work reviewed trades in equity markets. Among the critic-only studies, Sim and Kirk \cite{sim_generalized_2023} trained a Deep Q-Network on daily bars from 2006 to 2022. They tested it on assets it had not seen, and reported that the policy generalised to them. Théate and Ernst \cite{theate_application_2021} used the same algorithm on daily bars from 2012 to 2019. They reported consistent results in both calm and volatile periods. Their agent explored through an inverted copy of the environment, instead of through a noise process. Li et al. \cite{li_deep_2019} compared two agents on minute bars of stocks and futures from 2010 to 2020. The first was a Deep Q-Network over a denoising autoencoder. The second was Asynchronous Advantage Actor Critic over a long short-term memory network. Both beat the benchmark. The asynchronous actor-critic converged better, and the authors attribute this to its parallel training.

The actor-critic studies are closer to the design used in this work. Xiong et al. \cite{xiong_practical_2018} applied DDPG to daily bars from 2009 to 2018 and outperformed the benchmark. Yousefi \cite{yousefi_deep_2022} compared DDPG with Advantage Actor Critic on daily bars from 2009 to 2021. DDPG was the more stable of the two, and the study attributes this to its replay buffer and its noise process. Gran et al. \cite{gran_deep_2019} ran DDPG on monthly bars from 2005 to 2019. It still outperformed the benchmark after explicit and implicit costs were charged. Azhikodan et al. \cite{saini_stock_2019} added news sentiment to a state space that already held price data and technical indicators. They reported that this improved the results. However, the evaluation covers about five months, and the exact period is not given. Majidi et al. \cite{majidi_algorithmic_2022} used the Twin Delayed variant of DDPG on daily bars from 2010 to 2021, on both stocks and cryptocurrencies, with a continuous action in \([-1, 1]\). It beat the cryptocurrency benchmark but not the equity benchmark.

Sagiraju and Mogalla \cite{sagiraju_application_2021} need a separate mention. They compared four algorithms on daily bars from 1992 to 2020. All four beat the benchmark. Proximal Policy Optimisation and the Deep Q-Network were the more consistent, and DDPG and Advantage Actor Critic did worse than expected. It is the only study in the set that reports DDPG doing worse than the other algorithms, so it is a useful contrast to the studies above.

Two more studies combine algorithms instead of choosing one of them. Yang et al. \cite{yang_deep_2020} ran DDPG, Advantage Actor Critic and Proximal Policy Optimisation together on daily bars from 2009 to 2020. They switched between the three based on a turbulence measure. The ensemble beat the benchmark with explicit costs charged. Kong and So \cite{kong_empirical_2023} combined the same three algorithms on daily bars from 2016 to 2020. Their ensemble had unstable returns and underperformed the benchmark. Together, the two results indicate that combining agents does not by itself make the results robust. Wu et al. \cite{wu_adaptive_2020} sit between the two groups. They compared a gated recurrent Q-network with a gated recurrent policy-gradient agent on daily bars from 2008 to 2018, and the policy-gradient version did better.

\subsection{Other Market Implementations}

Three articles trade instruments other than equities and currencies. Chen and Gao \cite{chen_application_2019} traded exchange-traded funds on daily bars from 2000 to 2018. They compared a Deep Q-Network over a fully connected network with its recurrent version, and the recurrent model did better. Zhang et al. \cite{zhang_deep_2020} applied Advantage Actor Critic over a long short-term memory network to futures on daily bars from 2005 to 2019. It still outperformed the benchmark after costs. Jia et al. \cite{jia_lstm-ddpg_2021} also traded exchange-traded funds, on daily bars from 2002 to 2020, with DDPG over a long short-term memory network. They tested two versions. One held a fixed position size, and the other set the position size continuously. The version with a variable size did better, with explicit and implicit costs charged. This is the closest study in the set to the design used in this work. It is the only one that uses a continuous action to size a position and also charges for trading.

\subsection{Comparative Findings}

Reading the two tables column by column, instead of row by row, shows the choices that most studies share.

The state space is usually price data plus technical indicators computed from it. Studies that report strong results almost always include indicators, and not price alone \cite{carapuco_reinforcement_2018, huang_financial_2018, yang_deep_2020}. This agrees with the argument in \Cref{Section:ReinforcementLearning}. A raw price series is not Markovian, and indicators put enough of the recent history into the state for the property to hold approximately.

The action space is discrete in 15 of the 18 studies, usually with three choices: buy, hold and sell \cite{rundo_deep_2019, saini_stock_2019, xiong_practical_2018}. Three studies use a continuous action \cite{majidi_algorithmic_2022, jia_lstm-ddpg_2021, zhang_deep_2020}. Only one of them uses it to size a position, and not to give a direction \cite{jia_lstm-ddpg_2021}. A discrete action set is easier to learn. However, it takes position sizing out of the policy, so the size has to be set by a rule that the agent cannot adapt. The reward is the total return in 14 of the 18 studies. Three use the Sharpe ratio \cite{kong_empirical_2023, sagiraju_application_2021, jia_lstm-ddpg_2021} and one uses the log return \cite{majidi_algorithmic_2022}. A risk-adjusted reward costs more to compute during training, because it needs a window of past returns and not only the latest one. The studies that use it do not report a consistent gain from it.

For the bias-variance trade-off, the usual methods are feature engineering, normalisation, regularisation and changes to the network itself. Typical examples are denoising \cite{li_deep_2019}, and normalisation together with regularisation \cite{theate_application_2021}. For exploration, the \(\varepsilon\)-greedy policy is standard when the action set is discrete, and additive Gaussian noise when it is continuous \cite{wu_adaptive_2020, majidi_algorithmic_2022}.

Trading costs are not recorded in the same way across studies. Five of the 18 studies state that their headline result still holds with explicit and implicit costs \cite{huang_financial_2018, gran_deep_2019, yang_deep_2020, jia_lstm-ddpg_2021, zhang_deep_2020}. The others report returns without saying what was deducted, so the reader has to guess whether the figure is gross or net. Training usually uses a train-validation-test split, sometimes with cross-validation. One study trains online and keeps learning as new data arrives \cite{huang_financial_2018}. Of all the studies, this one comes closest to the conditions of live deployment.

The data covers long periods, but at a coarse resolution. Twelve of the 18 studies cover ten years or more, and the median span is eleven years. So a long history is normal in this field. The studies are also alike in resolution, which is low: 17 of the 18 work on bars, 13 of those on daily bars, and only one uses quote-level data \cite{carapuco_reinforcement_2018}. This matters for how costs are counted. A bar reports four prices and hides the order in which the price reached them.